\paragraph{Scoped formal support.}
The companion Lean development proves normalized finite-character Parseval
identities for an injective family of additive characters, preservation under
specified unit-modulus coefficient phases, and the translation identity on a
finite abelian group. It specializes to characters of $\mathbb Z/q\mathbb Z$
and proves injectivity of the centered residue modes under
$2K+1\leq q$. Separately, for a supplied complex representation of a finite
group, it proves the formula, idempotence, character-eigenvector range and
commutation properties of the character-weighted average, including the
permutation-sign specialization. Scalar lemmas prove the total error bound
from the four assumed intermediate energy inequalities, the fair-test
identity, quantitative YES/NO margins, and soundness after a subnormalized
weight and an assumed acceptance approximation error.

These are supporting finite-algebraic and scalar statements. They do not
formalize continuous torus integration or the production of phases by real
translations, the product-grid and spin realization, the concrete particle
representation, or Hilbert-space self-adjointness and orthogonality of the
average. Hardy/IMS localization, the exterior wedge, the Coulomb cap,
operator-domain and Ritz estimates, the Lipschitz quadrature argument,
quantum instruments and their polynomial circuit approximations, and
entangled-witness amplification remain in the human proof above. In
particular the assumed scalar energy and acceptance bounds are not Lean
proofs of their analytic or circuit hypotheses. Exact declarations,
compiler/dependency versions, source hashes, controls and coverage limits
are recorded in the accompanying support files; no complete formalization
of Theorem~\ref{thm:main} is claimed.
